\begin{abstract}
Sensor records of mechanical systems are seldom complete, and the missing dynamics must be inferred. The full--partial reconstruction mapping (FPRM) uses Takens' embedding theorem to learn the map between the delay reconstructions of a complete and an incomplete signal. The FPRM is extended and tested on the spring-coupled inverted pendulum (PIAR), a minimal model of inverted-pendulum seismic isolators controlled by one parameter $\mu=kL/(Mg)-1$. Five canonical regimes, from a libration inside the upright well to rotations near and far from the separatrix, are characterized with closed-form invariants and trajectories conserving energy to $10^{-10}\,MgL$. Under damping and forcing, a Feigenbaum cascade leads to a strange attractor with largest Lyapunov exponent $0.096\,(g/L)^{1/2}$ and Kaplan--Yorke dimension $2.28$. A cross-mapping analysis shows that the reflection $\theta\to-\theta$ makes the spring-force sign unrecoverable from the collar height unless the forcing phase is known. To represent such ambiguity, the Gaussian-process engine of the FPRM is replaced by a conditional neural spline flow, optionally regularized by equation-of-motion residuals. With the spring force observed and $90\%$ of the collar-height and angular-velocity records missing, the physics-regularized flow lowers the continuous ranked probability score 4.5-fold relative to the original FPRM (threefold when its Gaussian process uses all labeled samples) and restores the fractal Poincar\'e section; for ambiguous targets, the unregularized flow preserves both branches, whereas residuals blind to the branch can induce mode collapse. An open laboratory, comprising a browser-based interactive simulation and a Python package that regenerates every figure and number, makes the PIAR a transparent benchmark for dynamics-informed reconstruction.
\end{abstract}
